\section{Related Work}
\label{sec:rw}
\label{sec:rw-sec}

\paragraph{Agent runtimes and orchestration.}
AutoGen, LangGraph, and related systems provide abstractions for agent
coordination, tool invocation, and workflow state~\cite{wu2023autogen,langgraph2024}.
MCP standardizes communication between model-facing clients and tool/data
servers~\cite{anthropic2024mcp}. These projects are relevant integration
contexts, not equivalent security baselines: their primary goals differ
from multi-tenant resource and permission governance. Neuron can sit below
an orchestrator or beside a protocol gateway. Our comparison is therefore
architectural---which identity, pressure, and verdict signals must cross the
boundary---rather than a performance ranking against those projects.

AIOS-like work moves scheduling, memory, and tool management into an
agent-oriented systems layer~\cite{mei2025aios}. Neuron shares the premise
that agents require runtime support, but focuses on the interaction between
resource admission and permission availability. The key distinction is not
whether a framework has a rate limiter or an access-control list in
isolation; it is whether policy-plane capacity and timeout semantics remain
safe under tenant pressure and whether pressure can be part of a permission
decision when desired.

Recent agent-security surveys and benchmarks study policy compliance, tool
misuse, and interaction-level safety~\cite{wang2024policy,ruan2024toolemu,zhang2024agentsafetybench,struppek2024agent}. These works motivate runtime
enforcement, but do not evaluate whether tenant-induced executor contention
makes the authorization path unavailable or how timeout semantics convert
that availability loss into an admission decision.

\paragraph{Resource isolation and admission control.}
Operating systems and containers provide CPU, memory, process, and I/O
isolation. Token buckets and cgroups are established mechanisms for limiting
bursts and allocating capacity. A deployment can place its permission
service in a separate process or cgroup, which is a strong alternative to
Neuron's reserved executor. Our isolated baseline reflects the same design
principle inside one JVM. The paper does not claim that cross-layer coupling
replaces isolation: the evaluation shows that isolation protects
availability, whereas coupling supplies a pressure-conditioned verdict.

Work on overload control and circuit breaking similarly distinguishes
admission, queuing, and fallback behavior. Neuron's failure model applies
that distinction to a security-sensitive policy path. A fail-open fallback
trades security for availability; a fail-closed fallback makes the opposite
trade. Making both baselines explicit is essential because lower timeout
correctness in a fail-open system is otherwise easily misattributed to a new
security mechanism.

\paragraph{Access control, delegation, and confused deputies.}
Privilege attenuation and the confused-deputy problem have long motivated
capability systems and least-privilege delegation~\cite{hardy1988confused}.
Neuron's spawn-time privilege non-increase applies that principle to dynamic
sub-agent creation: a child cannot silently receive a more permissive
profile than its parent. Depth is an additional, agent-specific risk signal,
not a substitute for capabilities or target-scoped authorization. The
pressure-aware rule further treats runtime state as context for a narrow
class of destructive requests.

\paragraph{Policy availability and fail-safe defaults.}
Security engineering commonly recommends fail-safe defaults, but real
systems sometimes introduce permissive fallbacks to meet latency or
availability targets. Prior work on reference monitors emphasizes complete
mediation and tamper resistance; distributed authorization work also
recognizes policy-service latency and partitions. Our contribution is a
controlled characterization of this tradeoff in a multi-tenant agent
runtime, including an isolated fail-closed baseline. We do not claim the
general discovery that overload can delay a service. The contribution lies
in tracing that delay through agent permission semantics and evaluating a
resource-informed policy response.

Policy-as-code engines such as Open Policy Agent (OPA)~\cite{opa2024} provide
an important practical comparator: they separate declarative policy evaluation
from application code and can be deployed beside an API gateway or admission
controller. OPA, identity-oriented authorization, and relationship-based
authorization are complementary policy mechanisms; none by itself specifies
reserved executor capacity or whether a deadline miss under tenant pressure
fails open. Neuron evaluates that runtime availability and verdict boundary
explicitly rather than claiming to replace these policy engines.

\paragraph{Cross-layer observability.}
Distributed tracing correlates events across services through propagated
identifiers. Neuron uses the same systems idea across resource admission,
spawn, and permission paths. Correlation alone does not enforce safety; it
makes the causal chain inspectable and supplies consistent context to the
joint decision. The enforcement contribution is the policy input and its
reserved execution path, not the trace identifier by itself.

\paragraph{Positioning for software systems research.}
The work is a software-architecture study of a runtime control plane. Its
unit of comparison is an architecture configuration---shared fail-open,
shared fail-closed, isolated, or coupled---and its main outcomes are
policy-plane availability and verdict semantics. Sandboxing, WASM, Docker,
LLM routing, and formal models are useful components of the broader Neuron
prototype, but they are outside the causal evaluation and are not presented
as independent innovations here.
